\section{Quicksort}

\textbf{Idea:} Quicksort is a divide-and-conquer algorithm that sorts by
\textbf{partitioning} an array around a chosen pivot.
All elements smaller than the pivot move to the left,
and all elements greater stay on the right.
Then we recursively quicksort the left and right partitions. For simplicity, we choose the \textbf{last element as the pivot}.

% =========================================================
% Visualization (Partition) -- ALL steps
% =========================================================
\subsection{Visualization (Partition)}
Example: \([5,2,4,1,3]\) with pivot = last element (3)

% Usage:
% \QuickStep{a0}{a1}{a2}{a3}{a4}{leftIndex}{iIndex}{note}
% Set iIndex = -1 to hide i (used after loop), set leftIndex = -1 to hide left.
\newcommand{\QuickStep}[8]{%
\noindent\begin{tikzpicture}[x=1cm,y=1cm,>=Stealth]
  \def\n{5}
  \def\cellW{1.25}
  \def\H{0.85}

  % Outer box
  \draw[line width=1.2pt] (0,0) rectangle ({\n*\cellW},\H);

  % Dividers
  \foreach \k in {1,...,4} {
    \draw[line width=1.2pt] ({\k*\cellW},0) -- ({\k*\cellW},\H);
  }

  % Values
  \node at ({0.5*\cellW},0.5*\H) {#1};
  \node at ({1.5*\cellW},0.5*\H) {#2};
  \node at ({2.5*\cellW},0.5*\H) {#3};
  \node at ({3.5*\cellW},0.5*\H) {#4};
  \node at ({4.5*\cellW},0.5*\H) {#5};

  % Pivot pointer (always last index = 4)
  \draw[->,thick] ({(4+0.5)*\cellW},1.70) -- ({(4+0.5)*\cellW},\H);
  \node[anchor=south] at ({(4+0.5)*\cellW},1.70) {$pivot$};

  % If left and i overlap, draw one pointer labeled "left, i"
  \ifnum#6>-1
    \ifnum#7>-1
      \ifnum#6=#7
        \draw[->,thick] ({(#6+0.5)*\cellW},1.45) -- ({(#6+0.5)*\cellW},\H);
        \node[anchor=south] at ({(#6+0.5)*\cellW},1.45) {$left,\ i$};
      \else
        % left pointer (higher)
        \draw[->,thick] ({(#6+0.5)*\cellW},1.60) -- ({(#6+0.5)*\cellW},\H);
        \node[anchor=south] at ({(#6+0.5)*\cellW},1.60) {$left$};

        % i pointer (lower)
        \draw[->,thick] ({(#7+0.5)*\cellW},1.35) -- ({(#7+0.5)*\cellW},\H);
        \node[anchor=south] at ({(#7+0.5)*\cellW},1.35) {$i$};
      \fi
    \else
      % only left shown
      \draw[->,thick] ({(#6+0.5)*\cellW},1.45) -- ({(#6+0.5)*\cellW},\H);
      \node[anchor=south] at ({(#6+0.5)*\cellW},1.45) {$left$};
    \fi
  \else
    % left hidden, maybe i shown
    \ifnum#7>-1
      \draw[->,thick] ({(#7+0.5)*\cellW},1.45) -- ({(#7+0.5)*\cellW},\H);
      \node[anchor=south] at ({(#7+0.5)*\cellW},1.45) {$i$};
    \fi
  \fi

  % Note below
  \node[anchor=west] at (0,-0.45) {#8};
\end{tikzpicture}\par\vspace{0.35cm}
}

% -------------------------
% ALL partition steps
% -------------------------
\QuickStep{5}{2}{4}{1}{3}{0}{0}{Start: $left=0,\ i=0,\ pivot=3$}

\QuickStep{5}{2}{4}{1}{3}{0}{0}{$a[i]=5 > 3 \Rightarrow$ do nothing, $i \leftarrow i+1$}

\QuickStep{5}{2}{4}{1}{3}{0}{1}{$a[i]=2 < 3 \Rightarrow$ swap $a[left]$ and $a[i]$, $left \leftarrow left+1$}

\QuickStep{2}{5}{4}{1}{3}{1}{2}{After swap: $[2,5,4,1,3]$, now $left=1,\ i=2$}

\QuickStep{2}{5}{4}{1}{3}{1}{2}{$a[i]=4 > 3 \Rightarrow$ do nothing, $i \leftarrow i+1$}

\QuickStep{2}{5}{4}{1}{3}{1}{3}{$a[i]=1 < 3 \Rightarrow$ swap $a[left]$ and $a[i]$, $left \leftarrow left+1$}

\QuickStep{2}{1}{4}{5}{3}{2}{-1}{After swap: $[2,1,4,5,3]$, loop ends ($i=e$)}

\QuickStep{2}{1}{4}{5}{3}{2}{-1}{Final step: swap pivot with $a[left]$ (swap $3$ and $4$)}

\QuickStep{2}{1}{3}{5}{4}{2}{-1}{Done partition: pivot placed at index $left$ $\Rightarrow [2,1,3,5,4]$}

After placing the pivot in its correct position, quicksort then sorts the left subarray $[2,1]$ using $1$ as the pivot and the right subarray $[5, 4]$ using $4$ as the pivot, since the previous pivot $3$ is in its final spot. 
% =========================================================
% Code (match the video pseudocode)
% =========================================================
\subsection{Quicksort Code (Java)}
\begin{verbatim}
void quickSort(int[] arr, int s, int e) {
    if (e - s + 1 <= 1) return;

    int pivot = arr[e];
    int left = s;

    // partition
    for (int i = s; i < e; i++) {
        if (arr[i] < pivot) {
            int temp = arr[left];
            arr[left] = arr[i];
            arr[i] = temp;
            left++;
        }
    }

    // move pivot between the partitions
    int temp = arr[left];
    arr[left] = arr[e];
    arr[e] = temp;

    // recurse on left and right partitions
    quickSort(arr, s, left - 1);
    quickSort(arr, left + 1, e);
}
\end{verbatim}

\subsection{Time Complexity}
\begin{itemize}
  \item \textbf{Best case:} $O(n\log n)$ (pivot splits evenly)
  \item \textbf{Average case:} $O(n\log n)$
  \item \textbf{Worst case:} $O(n^2)$ (pivot gives very unbalanced splits)
\end{itemize}

\subsection{Space Complexity}
\begin{itemize}
  \item In-place partitioning uses $O(1)$ extra array space
  \item Recursion stack: $O(\log n)$ average, $O(n)$ worst
\end{itemize}